% Extension prepared by the Part F agent and independently reviewed by root.
% Do not infer the universal fixed-point assertion D5 from the example below.
\subsection{An abstract action outside every finite quotient}

\begin{claim}{D.10}{Proved}
Let $\mathcal U$ be the universal pro-cover of
$U=\mathbb P^1_{\overline\Q}\setminus\{0,1,\infty\}$ and
$\Omega$ its function field. There are a continuous tangential
semilinear action $\rho_0$ of $G_\Q$, a completed cusp $c$, and
an injective homomorphism
\[
 T:(\mathbb A_f,+)\longrightarrow\operatorname{Aut}_{\overline\Q}(\mathcal U)
\]
such that every $T_a$ fixes $c$,
\[
 \rho_0(g)T_a\rho_0(g)^{-1}=T_{\chi(g)a},
 \qquad
 \operatorname{Stab}_{\mathbb A_f}(\lambda)=\widehat\Z.
\]
Here $\lambda$ is the coordinate on $U$, viewed in $\Omega$,
and $\mathbb A_f$ denotes the additive finite adeles.
\end{claim}
\begin{proof}
Fix $\overline\Q\hookrightarrow\C$. The classical modular
uniformization identifies
$\mathfrak H/\Gamma(2)$ with $U(\C)$, and for
$q=\exp(\pi i\tau)$ its coordinate has expansion
\[
 \lambda(q)=16q-128q^2+704q^3+\cdots\in\Q[[q]].
\]
The product formula and this uniformization are given in
\citet[p.~628, equation~(1.0.2) and the following paragraph]{F_CDT2025}.
Only this classical coordinate description is used, not that
paper's unbounded-denominators theorem. Analytic finite-cover
comparison is SGA~1, XII, Theorem~5.1 \citep{F_SGA1}.
Choose the compatible cusp above $\tau=i\infty$ to embed
\[
 \Omega\hookrightarrow
 \mathcal P:=\bigcup_{n\ge1}\overline\Q((q^{1/n})).
\]
An element of $\Omega$ lifts to a meromorphic function on
$\mathfrak H$, invariant under a finite-index subgroup of
$\Gamma(2)$ and meromorphic at its cusps. Its selected Puiseux
germ has algebraic coefficients. The coordinate ring of
$\mathcal U$ consists of the elements with no poles on
$\mathfrak H$; this will also check the affine pro-curve, not
just its field.

Every finite adele is $a=r+z$ with $r\in\Q$ and
$z\in\widehat\Z$, and $\Q\cap\widehat\Z=\Z$.
Thus the quotient $\mathbb A_f/\widehat\Z\simeq\Q/\Z$
defines the character
\[
 \psi:\mathbb A_f\longrightarrow\mu_\infty(\C),
 \qquad \psi(r+z)=\exp(2\pi i r).
\]
For a Puiseux series set
\[
 T_a\left(\sum_k b_kq^{k/n}\right)
       =\sum_k b_k\psi(ka/n)q^{k/n}.
\]
Changing a denominator $n$ to $mn$ changes $k$ to $mk$,
so this formula is well-defined. It gives
$T_aT_b=T_{a+b}$, fixes constants, and preserves the
$q$-valuation.

If $z\in\widehat\Z$, then $T_z$ fixes $q$ and hence
$\lambda(q)$. It is the usual cusp-inertia action on the
Puiseux roots. Since $\Omega/\overline\Q(\lambda)$ is normal
with allowed branch locus $\{0,1,\infty\}$, this preserves
$\Omega$ and the affine coordinate ring.
For $r\in\Q$, the operator $T_r$ is analytically
$f(\tau)\mapsto f(\tau+2r)$. The matrix
$u(2r)=\left(\begin{smallmatrix}1&2r\\0&1\end{smallmatrix}\right)$
commensurates $\Gamma(2)$: clearing a denominator $d$ shows that
conjugation of $\Gamma(2d^2)$ lands in $\Gamma(2)$, and the
same argument applies to the inverse. Therefore if $f$ is
invariant under $\Gamma_f$ of finite index in $\Gamma(2)$,
then $T_rf$ is invariant under the finite-index subgroup
\[
 \Gamma(2)\cap u(2r)^{-1}\Gamma_f u(2r).
\]
It is meromorphic at every cusp. Consequently $T_rf$ is
algebraic over $\C(\lambda)$ with branching only over
$0,1,\infty$.

We check descent here explicitly. The Puiseux coefficients of
$T_rf$ remain algebraic. Take a polynomial relation over
$\C(\lambda)$ and clear denominators. Its finitely many
monomials in $\lambda$ and $T_rf$ have Laurent coefficients in
$\overline\Q$. The homogeneous linear system obtained by
comparing coefficients has a nonzero complex solution, hence
a nonzero solution over $\overline\Q$: its rank is determined
by finite minors. Thus $T_rf$ is algebraic over
$\overline\Q(\lambda)$ and belongs to $\Omega$.
Moreover $\tau\mapsto\tau+2r$ preserves $\mathfrak H$ and
the rational cusps, so it preserves absence of poles on
$\mathfrak H$. The inverse is $T_{-r}$. Every $a=r+z$ now
acts on $\mathcal U$ as asserted. This proof does not assume
uniqueness of roots in the full automorphism group.

These automorphisms also act on the completed pro-cover.
Indeed the image of any finite-level function field is
contained in a further finite level. The resulting rational
map between the smooth proper compactifications extends
uniquely everywhere, by properness and normality for curves.
The extensions are compatible under composition and passage
to further levels. Applying the inverse operator proves that
they define an automorphism of the completed pro-object.

The coefficient of $q$ in $T_a\lambda$ is $16\psi(a)$.
It follows that $T_a\lambda=\lambda$ if and only if
$a\in\widehat\Z$: sufficiency follows because $T_a$ then
fixes $q$. If such an $a$ acts trivially on all of $\Omega$,
apply it to $\lambda^{1/n}\in\Omega$ for every $n$.
Their leading terms force $\psi(a/n)=1$, so
$a\in\bigcap_n n\widehat\Z=0$. This proves faithfulness.

Let $\rho_0(g)$ act coefficientwise on $\mathcal P$, fixing
every $q^{1/n}$. It restricts to the continuous tangential
action on $\Omega$ supplied by
$\operatorname{Spec}\Q((q))\to U$, $\lambda\mapsto\lambda(q)$.
The parameter $q$ is a rational formal uniformizer; it need
not belong to $\Omega$. Directly on each Puiseux monomial,
\[
 g\bigl(\psi(a/n)\bigr)=\psi\bigl(\chi(g)a/n\bigr)
\]
proves the displayed conjugation identity. Both types of
operator preserve the selected valuation, hence its completed
cusp $c$.
\end{proof}

\begin{claim}{D.11}{Proved}
There is an abstract semilinear action $\rho$ of $G_\Q$ on
$\mathcal U$ that normalizes no deck group of a finite
hyperbolic quotient. Its $G_\Q^\circ$-fixed set on the
completed pro-cover is precisely the single cusp $c$ of D.10.
This answers the second demand of D5. It does not settle
D5's universal fixed-point assertion.
\end{claim}
\begin{proof}
Fix an odd prime $\ell$ and distinct rational primes
$p_1,p_2,\ldots$ different from $\ell$. Compatible radicals
give Kummer cocycles and their continuous product
\[
 \kappa:G_\Q\longrightarrow V:=\prod_{i\ge1}\Z_\ell,
 \qquad
 \kappa(gh)=\kappa(g)+\chi_\ell(g)\kappa(h).
\]
We have chosen a compatible basis of the Tate twist in
writing these coordinates. Its restriction to $G_\Q^\circ$
is a continuous surjective homomorphism. Indeed, for finitely
many $p_i$ and $m\ge1$, put
\[
 K=\Q(\mu_{\ell^m}),\qquad
 L=K(p_i^{1/\ell^m}).
\]
The extension $K/\Q$ is unramified at the $p_i$.
Valuations there show that their classes in
$K^\times/(K^\times)^{\ell^m}$ are independent of order
$\ell^m$. Kummer theory gives
$T:=\operatorname{Gal}(L/K)\simeq(\Z/\ell^m)^r$.
The extension $L/\Q$ is Galois, and a lift of a cyclotomic
automorphism with character $2$ modulo $\ell$ acts on $T$
by that scalar. Its commutators with $T$ exhaust $T$, since
$\chi_\ell-1$ is a unit. As the quotient by $T$ is abelian,
$T$ is exactly the commutator subgroup. Hence
$L\cap\Q^{\mathrm{ab}}=K$, which proves the required
surjection from $G_\Q^\circ$ at every finite level.
Compactness gives the surjection onto $V$.

Let $e_i\in V$ be the coordinate vectors. They are linearly
independent in $V\otimes_{\Z_\ell}\Q_\ell$. Extend them to
a vector-space basis and choose a linear functional taking
$e_i$ to $\ell^{-i}$. Its restriction is a
$\Z_\ell$-linear map
\[
             f:V\longrightarrow\Q_\ell,
                   \qquad f(e_i)=\ell^{-i}.
\]
This basis choice uses the axiom of choice and specifies an
abstract cocycle; it is not an effective algorithm for
evaluating the cocycle on arbitrary Galois elements.
It is onto, because its image contains
$\ell^{-i}\Z_\ell$ for every $i$. Its kernel is dense:
realize any finitely many prescribed coordinates by a
finitely supported $v$. Choose $j$ outside those coordinates
so large that $\ell^j f(v)\in\Z_\ell$, and use
$v-\ell^j f(v)e_j$. This has the prescribed coordinates
and is killed by $f$.

Regard $\Q_\ell$ as the subgroup of $\mathbb A_f$ supported
at $\ell$ and set $a_g=f(\kappa(g))$. Linearity gives
$a_{gh}=a_g+\chi(g)a_h$. Therefore
\[
                   \rho(g)=T_{a_g}\rho_0(g)
\]
is an action by D.10, semilinear for the specified usual
action on $\overline\Q$. Moreover
\[
 a(G_\Q^\circ)=\Q_\ell,
 \qquad K_0:=\ker(a|_{G_\Q^\circ})
                     \text{ is dense in }G_\Q^\circ.
\]
For density, a continuous surjective homomorphism of compact
groups is an open quotient map; apply this to $\kappa$ and
the dense kernel of $f$.

The orbit of $\lambda$ is infinite: $\rho_0$ fixes this
rational coordinate and D.10 gives
\[
 \{\rho(g)\lambda:g\in G_\Q^\circ\}
       \simeq \Q_\ell/\Z_\ell.
\]
We show why this rules out \emph{every} normalized finite
hyperbolic quotient, without using a congruence-quotient
comparison. If a deck group $H$ were normalized, the action
would descend to its hyperbolic quotient $W$ and the field
$F=\Omega^H$. Choose a number field $K$ defining $W$ and
finitely many function-field generators. For $g\in G_K$,
compare the descended semilinear automorphism to the
coefficient action of that $K$-model. Their quotient is in
the finite group $\operatorname{Aut}_{\overline\Q}(W)$.
Thus those generators have finite $G_K$-orbits, and then
finite $G_\Q$-orbits. Constants also have finite orbits, so
every element of $F$ does. Each element of $\Omega$ is
algebraic over $F$. Its orbit is among the roots of the
finitely many conjugates of a polynomial equation over
$F$, hence is finite. This contradicts the orbit of
$\lambda$. No continuity was presumed in this argument.

It remains to determine the fixed set. Every $T_a$ and
$\rho_0(g)$ fixes $c$, so $c$ is fixed by $\rho$.
If a completed geometric closed pro-point $P$ is fixed by
$\rho(G_\Q^\circ)$, it is fixed by $\rho_0(K_0)$.
Its stabilizer under $\rho_0$ is closed: take cofinal
characteristic finite levels, descending over $\Q$ under
$\rho_0$. At each level an algebraic closed point has an
open and closed Galois stabilizer, and the stabilizer of
$P$ is their intersection. Density of $K_0$ implies that
$P$ is fixed by $\rho_0(G_\Q^\circ)$.

The same Kummer argument as in D.7 proves that the latter
fixed set is $\{c\}$. Explicitly, if $P\ne c$, choose a
characteristic covering level separating their projections
and dominating the degree-nine Fermat cover
$x^3+y^3=1$, $\lambda=x^3$, of $U$. Its smooth proper
compactification has positive genus. More explicitly, if
$N\subset\Pi_U$ is this characteristic subgroup, then
$N\rho_0(G_\Q)\subset\pi_1(U_\Q)$ defines a finite
\'etale curve over $\Q$; take its smooth proper
compactification over $\Q$. The two projected
$\Q^{\mathrm{ab}}$-points have the same proper section
class, since the common action fixes both pro-points.
Their Abel--Jacobi difference has zero Kummer class and
therefore vanishes by F.3, based on the precise abelian
Kummer theorem of \citet[Theorem~3.1]{F_Tsu2020}.
Distinct points on a positive-genus smooth proper curve
cannot have principal difference: that would give a
degree-one map to $\mathbb P^1$. The contradiction proves
$P=c$.
\end{proof}

The constructed action lies outside the finite-quotient
normalization cases covered by R19, while satisfying the
predicted fixed-point conclusion itself. Thus automatic
finite-quotient descent cannot reduce all abstract actions
to B. This construction alone neither proves nor refutes the universal
assertion in B$'$, and does not exclude every logically different
reduction to B. The separate construction in D.12--D.14 below refutes
the universal assertion by an action with empty fixed set.
